\section{Introduction}\label{sec:introduction}

A polynomial threshold function assigns a sign to each vertex of the
Boolean cube by evaluating a real polynomial. On the cube, the identities
$x_i^2=1$ allow every polynomial to be replaced by a multilinear polynomial
of no larger degree without changing its values. Average sensitivity
measures the expected number of coordinate changes that reverse that sign.
For $x\in\{-1,1\}^n$, let $x^{\oplus i}$ be obtained by reversing
coordinate $i$. For $f:\{-1,1\}^n\to\{-1,1\}$, define
\begin{equation}\label{eq:influence-definition}
 I(f)=\sum_{i=1}^n\Prb\{f(X)\ne f(X^{\oplus i})\},
 \qquad X\text{ uniform on }\{-1,1\}^n.
\end{equation}
This quantity is also called the total influence of $f$. Throughout,
\[
 \sgn(t)=\begin{cases}1,&t\ge0,\\-1,&t<0.\end{cases}
\]
We prove the following estimate.

\begin{theorem}\label{thm:main}
Let $n\ge1$ and $1\le d\le n$ be integers. Let $p$ be a real multilinear
polynomial of degree at most $d$, and define
$f:\{-1,1\}^n\to\{-1,1\}$ by $f(x)=\sgn(p(x))$, with $\sgn(0)=1$.
Then, for average sensitivity under the uniform law,
\[
 I(f)\le 8d\sqrt n.
\]
\end{theorem}

The constant is independent of both $d$ and $n$. In particular, the degree
may grow with the dimension. The polynomial may vanish on the cube, and
no regularity condition is imposed on $p$.

Gotsman and Linial proposed an exact extremal bound for average
sensitivity~\cite[Section~5, p.~47]{GotsmanLinial1994}.
Their candidate is symmetric: it is the sign of a degree-$d$ polynomial
in $x_1+\cdots+x_n$ whose roots lie between the central levels of the
cube. The proposal asserts that this function maximizes influence among
all degree-$d$ polynomial threshold functions in $n$ variables.
Chapman constructed counterexamples to this exact assertion and
explicitly distinguished it from the asymptotic bound
$I(f)=O(d\sqrt n)$
\cite[Conjectures~1.1--1.2 and Theorem~1.1]{Chapman2018}.
Kim, Maldonado, and Wellens independently constructed quadratic
counterexamples for odd $n\ge5$ and proved influence bounds for
quadratic polynomials with restricted support
graphs~\cite[Theorems~1.1--1.3]{KimMaldonadoWellens2017}.
Theorem~\ref{thm:main} establishes the asymptotic bound with an absolute
constant; it makes no assertion about exact maximizers.
Numbered locators in references with a listed preprint refer to that
preprint version.

The classical symmetric examples explain the scale of the theorem.
Let $1\le d\le\sqrt n$, let $J$ consist of the $d$ indices nearest
$(n-1)/2$ in $\{0,\ldots,n-1\}$, breaking a tie arbitrarily, and put
\[
 t(x)=\frac{n+x_1+\cdots+x_n}{2},\qquad
 p_J(x)=\prod_{j\in J}\bigl(t(x)-j-\tfrac12\bigr).
\]
Here $t(x)$ is the number of coordinates of $x$ equal to $+1$.
Multilinearization preserves the values of $p_J$ on the cube and has
degree at most $d$. Its sign changes exactly across the boundaries
between Hamming weights $j$ and $j+1$ for $j\in J$. There are
$(n-j)\binom nj=n\binom{n-1}j$ unoriented edges at such a boundary,
so the ordered-edge normalization in~\eqref{eq:influence-definition}
gives
\[
 I(\sgn p_J)=n\,2^{-(n-1)}\sum_{j\in J}\binom{n-1}j
 \asymp d\sqrt n.
\]
Here the comparison constants are absolute. The central-binomial
estimate from Stirling's formula, together with the ratios of consecutive
binomial coefficients, shows that the binomial masses are comparable
to $n^{-1/2}$ within $\sqrt n/2$ of the center, which includes all
selected indices; the case $n=1$ follows directly from the count.
Thus the order $d\sqrt n$ is sharp throughout
this degree range; see also
\cite[p.~3, discussion following Theorem~1.1]{KimMaldonadoWellens2017}.
For every Boolean function, the additional bound $I(f)\le n$ holds.

Early general bounds established sublinear influence for every fixed
degree. For $n>1$, Harsha, Klivans, and Meka proved
$I(f)\le2^{O(d)}n^{1-1/(4d+6)}$
\cite[Theorem~1.6]{HarshaKlivansMeka2014}, while the independent work of
Diakonikolas, Raghavendra, Servedio, and Tan gave
$I(f)\le2^{O(d)}(\log n)n^{1-1/(4d+2)}$
\cite[Theorem~1.1]{DRST2014}.
These results first circulated in 2009, before their 2014 journal
publications. Their proofs use regularity and invariance arguments to
compare sufficiently regular Boolean polynomials with Gaussian ones.
Kane subsequently established the square-root exponent in the dimension:
\begin{equation}\label{eq:kane-bound}
 I(f)\le\sqrt n\,(\log n)^{O(d\log d)}2^{O(d^2\log d)}
 \qquad(n>1)
\end{equation}
\cite[Theorem~2]{Kane2014}.
For each fixed $d\ge2$, this is an $n^{1/2+o(1)}$ bound with a
polylogarithmic loss. His proof combines random restrictions, regularity,
and invariance estimates. Theorem~\ref{thm:main} removes that loss and
gives linear dependence on $d$ with an absolute constant.

Small moments of the local sensitivity
$s_f(x)=|\{i:f(x)\ne f(x^{\oplus i})\}|$ give further information
about its distribution. Chang, Slote, Volberg, and Zhang bounded the
Boolean surface area $\E\sqrt{s_f(X)}$ by a polynomial in $\log(en)$
for each fixed degree
\cite[Theorem~1.1, equation~(1.6)]{ChangSloteVolbergZhang2026}.
Tseng and Volberg studied a wider range of sensitivity moments
\cite{TsengVolberg2026}.
For comparison with the first moment, the pointwise inequality
$s_f\le\sqrt n\sqrt{s_f}$ converts the surface-area estimate into an
influence bound with a logarithmic loss. At moment order one, the
estimate in~\cite[Section~2.5, equation~(11)]{TsengVolberg2026}
is $C(d)\sqrt n(\log n)^{Cd\log d}$ for $n>1$, which likewise retains a
logarithmic factor for $d\ge2$.

\paragraph{Proof strategy.}
The proof separates a probability estimate from a finite-dimensional
operator construction. The probability estimate says that equal-law
random variables of variance $\sigma^2$ satisfy
$\E(U-V)^2\le8a\sigma$ whenever $V-U\le a$ almost surely.
A bound on increments in the opposite direction is not assumed.
Equality of the marginal laws supplies the needed balance through an
increasing function whose increments dominate squared displacement.
Section~\ref{sec:coupling} proves this estimate directly.

For the operator construction, begin with the spaces of cube polynomials
of degrees at most $k$, for $0\le k\le n$. Multiply these spaces by
$\sqrt w$ for a positive weight $w$, and take their successive orthogonal
increments in the counting inner product. The increments have dimensions
$\binom nk$, although the weight changes their positions. Give the
$k$th increment eigenvalue $k-n/2$; this defines a self-adjoint operator
$M$ with fixed squared Hilbert--Schmidt norm $n2^n/4$.
Its squared Hilbert--Schmidt norm is the sum of the squared moduli of
all matrix entries.
Multiplication by a coordinate connects only equal or adjacent grades,
because it raises polynomial degree by at most one and is self-adjoint.
These are the two features of the grading used in the proof: fixed
dimensions and locality under coordinate multiplication
(Section~\ref{sec:grading}).

When $w\equiv1$, $M$ is minus one-half of the cube adjacency matrix.
For a general weight, some edge entries may have smaller squared modulus,
and other entries may appear. Section~\ref{sec:edges} shows that the total squared
entry mass between vertices at distance at least two is bounded by the
mass on the diagonal.
Together with a bound on each edge entry, this controls the total deficit
from the unweighted squared edge value $1/4$.
If $H$ multiplies by a sign function $h$, the anticommutator $MH+HM$
retains diagonal entries and equal-sign edge entries, multiplying them by
$2$ or $-2$. Its squared norm therefore bounds the number of equal-sign
edges even for an arbitrary positive weight.

Finally, multiply $f$ by full parity to obtain
$h(x)=(\prod_{i=1}^n x_i)f(x)$, and use $w=|p|$ after removing zeros
without changing signs. The equal-sign edges of $h$ are exactly the
sensitive edges of $f$. Moreover, multiplication by parity turns the
Fourier support of $p$ into degrees at least $n-d$.
This forces the block of $H$ between grades $r$ and $s$ to vanish when
$r+s<n-d$. The normalized squared block norms define a probability law
for $(R,S)$ whose two marginals are $\operatorname{Bin}(n,1/2)$.
Reflecting one index gives $U=S$ and $V=n-R$, with the same marginals and
$V-U\le d$. The coupling estimate then bounds the anticommutator norm,
and the edge estimate gives the theorem. Section~\ref{sec:parity}
performs this assembly, keeping track of ordered edges.

\paragraph{Consequences.}
Section~\ref{sec:consequences} derives a dimension-free Boolean noise
sensitivity bound and a quantitative learning guarantee from the influence
estimate. The learning statement concerns independent labeled samples
whose input marginal is uniform, while the labels may be arbitrary.
It follows by combining noise smoothing with polynomial regression.
